\subsection{SWAT\textsuperscript{+} project assembly and workflow outputs}
\label{sec:methods-assembly}

\begin{table}[htbp]
  \centering
  \caption{Representative SWAT\textsuperscript{+} output file classes.}
  \label{tab:package-manifest-abridged}
  \small
  \begin{tabularx}{\textwidth}{@{}lYl@{}}
\toprule
Output class & Purpose & Always? \\
\midrule
Vector shapefile & Channel geometry for QSWAT+ / SWAT+ GIS import & always \\
Vector shapefile & Reach watershed polygons & always \\
Vector shapefile & Dissolved subbasin polygons after one-outlet partition & always \\
Vector shapefile & Lake polygons when present & conditional \\
Weather text & PRISM-derived precipitation and temperature station files & always \\
JSON sidecar & Channel/subbasin processing statistics for QA & always \\
SWAT+ SQLite project & SWAT+ Editor / SQLite project tables after import pipeline & always \\
Project archive & Exported SWAT+ project archive bundle & conditional \\
\bottomrule
\end{tabularx}
%
\end{table}

For a declared vector processing unit (VPU) and list of HUC12 units (or the equivalent extent key in the implementation), the extractor clips the preprocessed stream network, applies the one-outlet subbasin logic and lake linkage rules of Section~\ref{sec:methods-nhd}, and writes vector products as ESRI shapefiles in the generated SWAT\textsuperscript{+} project tree intended for QSWAT\textsuperscript{+}-oriented workflows \citep{QSWATplusSoftware}.
Table~\ref{tab:package-manifest-abridged} summarizes primary output file classes; representative path patterns and integrity expectations are listed in Supplementary Table~S1.
Every manuscript benchmark domain then runs the standard modified-input generation post-pass (Section~\ref{sec:methods-nhd}): subbasin merge by the documented $A_{\min}$ rule, outlet-topology enforcement, and the separate tiny-basin guard before shapefiles are rewritten for QSWAT\textsuperscript{+} handoff.

Weather text inputs, SWAT\textsuperscript{+} text configuration, and metadata that accompany the vector exports are treated as part of the same generated \emph{project tree}.
The archived project manifest documents which additional files accompany the vector exports for integrity checks on the evaluation-protocol basins.
This contribution does not center on groundwater flow coupling, MODFLOW, or regional recharge workflows; those capabilities, where present in the broader implementation, are out of scope for the hydrography-first narrative developed here.
